\def\definitionsyntaxcolouring#1{\pdfliteral direct{0.27 0.27 0.27 rg}}
    \def\functionsyntaxcolouring#1{\pdfliteral direct{0.75 0 0 rg}}
    \def\identifiersyntaxcolouring#1{\pdfliteral direct{0.25 0.25 1 rg}}
    \def\elementsyntaxcolouring#1{\pdfliteral direct{0.25 0.25 0.5 rg}}
    \def\reservedsyntaxcolouring#1{\pdfliteral direct{0.375 0 0 rg}}
    \def\stringsyntaxcolouring#1{\pdfliteral direct{0.25 0.5 0.25 rg}}
    \def\charactersyntaxcolouring#1{\pdfliteral direct{0.125 0.25 0.125 rg}}
    \def\typesyntaxcolouring#1{\pdfliteral direct{0.125 0.5 0.125 rg}}
    \def\constantsyntaxcolouring#1{\pdfliteral direct{0.5 0.25 0.125 rg}}
    \def\plainsyntaxcolouring#1{\pdfliteral direct{0 0 0 rg}}
    \def\extractsyntaxcolouring#1{\pdfliteral direct{0.27 0.27 0.27 rg}}
    \def\commentsyntaxcolouring#1{\pdfliteral direct{0.25 0.25 0.25 rg}}
    \def\endcolouring#1{\special{PDF:0 g}}

    \def\inwebimage#1{\pdfximage{#1}\smallskip\noindent
    \hbox to\hsize{\hfill\pdfrefximage \pdflastximage\hfill}\smallskip}
    %
    \def\inwebimagewidth#1#2{\pdfximage width #2cm {#1}\smallskip\noindent
    \hbox to\hsize{\hfill\pdfrefximage \pdflastximage\hfill}\smallskip}
    %
    \def\inwebimageheight#1#2{\pdfximage height #2cm {#1}\smallskip\noindent
    \hbox to\hsize{\hfill\pdfrefximage \pdflastximage\hfill}\smallskip}
    %
    \def\inwebimagewidthheight#1#2#3{\pdfximage width #3cm height #2cm {#1}\smallskip\noindent
    \hbox to\hsize{\hfill\pdfrefximage \pdflastximage\hfill}\smallskip}
    %
    \def\inwebstartlinkurl#1{{\leavevmode\pdfstartlink
    attr{/Border [2 1 1] /C [0.5 0 0]}
    user{/Subtype /Link
    /A << /S /URI /URI (#1) >>}%
    }}
    %
    \def\inwebstartlinkpara#1{{\leavevmode\pdfstartlink
    attr{/Border [0 0 0] /C [0.9 0 0]}
    goto name {#1}}}%
    %
    \def\inwebstartlinkgoto#1{{\leavevmode\pdfstartlink
    attr{/Border [2 1 1] /C [0.5 0 0]}
    goto name {#1}}}%
    %
    \def\inweblinktourl#1#2{\inwebstartlinkurl{#1}{#2}\pdfendlink}
    \def\inweblinktopara#1#2{\inwebstartlinkpara{#1}{#2}\pdfendlink}
    \def\inweblinktoanchor#1#2{\inwebstartlinkgoto{#1}{#2}\pdfendlink}
    %
    \def\inwebanchor#1{\pdfdest name {#1} xyz }
    %
    % Beginning of `inweb-macros.tex`
    %
    % These are TeX macros for use by Inweb when rendering to a Plain TeX variant
    % (such as TeX itself, pdftex or xetex): they should work on any TeX. Macros
    % which do not have the luxury of engine-independence have been hived off into
    % the files `links-macros.tex`, `colours-macros.tex`, and `images-macros.tex`.
    %
    % (1) Fonts
    % There are more fonts here than we actually need for Inweb purposes, but it
    % seems easiest to follow Knuth's semi-standard book design code.
    %
    % See Appendix B of the TeXBook: some `\font` definitions, such as `\tenrm`
    % and `\tentt`, are already defined in Plain TeX.
    %
    \font\defnannotationfont=cmti9
    \def\defnannotationtext#1{{\defnannotationfont #1}}
    \font\tentex=cmtex10
    \font\sourcefont=cmss10
    \font\holonfont=cmss9
    \def\holonnametext#1{{\holonfont #1}}
    \hyphenchar\holonfont=-1
    \font\usagefont=cmss8
    \def\endnotetext#1{{\usagefont #1}}
    \def\sectiontoctext#1{{\usagefont #1}}
    \font\inchhighfont=cminch
    \font\reducedinchhighfont=cminch at 60pt
    \font\chaptertitlefont=cmss10 at 24pt
    \def\chaptertitletext#1{\chaptertitlefont #1}
    \font\sectiontitlefont=cmss10 at 20pt
    \def\sectiontitletext#1{\sectiontitlefont #1}
    \font\ssectiontitlefont=cmss10 at 20pt
    \font\sssectiontitlefont=cmr10 at 12pt
    \font\ninerm=cmr9
    \font\eightrm=cmr8
    \font\sixrm=cmr6
    \font\ninei=cmmi9
    \skewchar\ninei='177
    \font\eighti=cmmi8
    \skewchar\eighti='177
    \font\sixi=cmmi6
    \skewchar\sixi='177
    \font\ninesy=cmsy9
    \skewchar\ninesy='60
    \font\eightsy=cmsy8
    \skewchar\eightsy='60
    \font\sixsy=cmsy6
    \skewchar\sixsy='60
    \font\holonnumberfont=cmss10 at 7pt
    \def\holonnumbertext#1{{\holonnumberfont #1}}
    \font\eightss=cmssq8
    \font\eightssi=cmssqi8
    \font\ninebf=cmbx9
    \font\eightbf=cmbx8
    \font\sixbf=cmbx6
    \font\ninett=cmtt9
    \hyphenchar\ninett=-1
    \font\eighttt=cmtt8
    \hyphenchar\eighttt=-1
    \hyphenchar\tentt=-1 % inhibit hyphenation in typewriter type
    \font\ninesl=cmsl9
    \font\eightsl=cmsl8
    \font\nineit=cmti9
    \font\eightit=cmti8
    \font\tenu=cmu10 % unslanted text italic
    \font\magnifiedfiverm=cmr5 at 10pt
    \font\manual=manfnt % font used for the METAFONT logo, etc.
    \font\cmman=cmman % font used for miscellaneous Computer Modern variations
    %
    % Knuth's standard way to change point size:
    %
    \newskip\ttglue
    \def\tenpoint{\def\rm{\fam0\tenrm}%
      \textfont0=\tenrm \scriptfont0=\sevenrm \scriptscriptfont0=\fiverm
      \textfont1=\teni \scriptfont1=\seveni \scriptscriptfont1=\fivei
      \textfont2=\tensy \scriptfont2=\sevensy \scriptscriptfont2=\fivesy
      \textfont3=\tenex \scriptfont3=\tenex \scriptscriptfont3=\tenex
      \def\it{\fam\itfam\tenit}%
      \textfont\itfam=\tenit
      \def\sl{\fam\slfam\tensl}%
      \textfont\slfam=\tensl
      \def\bf{\fam\bffam\tenbf}%
      \textfont\bffam=\tenbf \scriptfont\bffam=\sevenbf
       \scriptscriptfont\bffam=\fivebf
      \def\tt{\fam\ttfam\tentt}%
      \textfont\ttfam=\tentt
      \tt \ttglue=.5em plus.25em minus.15em
      \normalbaselineskip=12pt
      \def\MF{{\manual META}\-{\manual FONT}}%
      \let\sc=\eightrm
      \let\big=\tenbig
      \setbox\strutbox=\hbox{\vrule height8.5pt depth3.5pt width\z@}%
      \normalbaselines\rm}
    %
    \def\ninepoint{\def\rm{\fam0\ninerm}%
      \textfont0=\ninerm \scriptfont0=\sixrm \scriptscriptfont0=\fiverm
      \textfont1=\ninei \scriptfont1=\sixi \scriptscriptfont1=\fivei
      \textfont2=\ninesy \scriptfont2=\sixsy \scriptscriptfont2=\fivesy
      \textfont3=\tenex \scriptfont3=\tenex \scriptscriptfont3=\tenex
      \def\it{\fam\itfam\nineit}%
      \textfont\itfam=\nineit
      \def\sl{\fam\slfam\ninesl}%
      \textfont\slfam=\ninesl
      \def\bf{\fam\bffam\ninebf}%
      \textfont\bffam=\ninebf \scriptfont\bffam=\sixbf
       \scriptscriptfont\bffam=\fivebf
      \def\tt{\fam\ttfam\ninett}%
      \textfont\ttfam=\ninett
      \tt \ttglue=.5em plus.25em minus.15em
      \normalbaselineskip=11pt
      \def\MF{{\manual hijk}\-{\manual lmnj}}%
      \let\sc=\sevenrm
      \let\big=\ninebig
      \setbox\strutbox=\hbox{\vrule height8pt depth3pt width\z@}%
      \normalbaselines\rm}
    %
    \def\ttninepoint{\def\rm{\fam0\ninerm}%
      \textfont0=\ninerm \scriptfont0=\sixrm \scriptscriptfont0=\fiverm
      \textfont1=\ninei \scriptfont1=\sixi \scriptscriptfont1=\fivei
      \textfont2=\ninesy \scriptfont2=\sixsy \scriptscriptfont2=\fivesy
      \textfont3=\tenex \scriptfont3=\tenex \scriptscriptfont3=\tenex
      \def\it{\fam\itfam\nineit}%
      \textfont\itfam=\nineit
      \def\sl{\fam\slfam\ninesl}%
      \textfont\slfam=\ninesl
      \def\bf{\fam\bffam\ninebf}%
      \textfont\bffam=\ninebf \scriptfont\bffam=\sixbf
       \scriptscriptfont\bffam=\fivebf
      \def\tt{\fam\ttfam\ninett}%
      \textfont\ttfam=\ninett
      \normalbaselineskip=11pt
      \normalbaselines\rm}
    %
    \def\eightpoint{\def\rm{\fam0\eightrm}%
      \textfont0=\eightrm \scriptfont0=\sixrm \scriptscriptfont0=\fiverm
      \textfont1=\eighti \scriptfont1=\sixi \scriptscriptfont1=\fivei
      \textfont2=\eightsy \scriptfont2=\sixsy \scriptscriptfont2=\fivesy
      \textfont3=\tenex \scriptfont3=\tenex \scriptscriptfont3=\tenex
      \def\it{\fam\itfam\eightit}%
      \textfont\itfam=\eightit
      \def\sl{\fam\slfam\eightsl}%
      \textfont\slfam=\eightsl
      \def\bf{\fam\bffam\eightbf}%
      \textfont\bffam=\eightbf \scriptfont\bffam=\sixbf
       \scriptscriptfont\bffam=\fivebf
      \def\tt{\fam\ttfam\eighttt}%
      \textfont\ttfam=\eighttt
      \tt \ttglue=.5em plus.25em minus.15em
      \normalbaselineskip=9pt
      \def\MF{{\manual opqr}\-{\manual stuq}}%
      \let\sc=\sixrm
      \let\big=\eightbig
      \setbox\strutbox=\hbox{\vrule height7pt depth2pt width\z@}%
      \normalbaselines\rm}
    %
    \def\tenmath{\tenpoint\fam-1 } % use after $ in ninepoint sections
    \def\tenbig#1{{\hbox{$\left#1\vbox to8.5pt{}\right.\n@space$}}}
    \def\ninebig#1{{\hbox{$\textfont0=\tenrm\textfont2=\tensy
      \left#1\vbox to7.25pt{}\right.\n@space$}}}
    \def\eightbig#1{{\hbox{$\textfont0=\ninerm\textfont2=\ninesy
      \left#1\vbox to6.5pt{}\right.\n@space$}}}
    %
    % (2) Inweb headings
    % `\weavesection` and its heavier friends are Inweb sections, the grander the
    % more `s` suffixes they have: untitled paragraph, titled paragraph, section, chapter.
    %
    \def\weavesectionsss#1#2#3#4#5{\vfill\eject\noindent\chaptertitletext{#2}\mark{#3}\vskip 0.5in\par\noindent\titlepage}
    \def\weavesectionss#1#2#3#4#5{\vfill\eject\noindent\sectiontitletext{#2\hfill #5}\mark{#3}\bigskip\par\noindent\titlepage}
    \def\weavesections#1#2#3#4#5{\bigskip\goodbreak\noindent{\bf #4#1. #2.}\mark{#3}\quad}
    \def\weavesection#1#2#3#4#5{\bigskip\goodbreak\noindent{\bf #4#1.}\mark{#3}\quad}
    %
    % The `t` variants are for themed extracts.
    \def\tweavesections#1#2#3#4#5{\bigskip\goodbreak\noindent{\bf #5.#4#1. #2.}\quad}
    \def\tweavesection#1#2#3#4#5{\bigskip\goodbreak\noindent{\bf #5.#4#1.}\quad}
    %
    % The `ns` variants have less vertical spacing before them.
    \def\nsweavesections#1#2#3#4#5{\noindent{\bf #4#1. #2.}\mark{#3}\quad}
    \def\nsweavesection#1#2#3#4#5{\noindent{\bf #4#1.}\mark{#3}\quad}
    %
    % (3) Strikethrough
    % Plain TeX does not support text struck through, out of the box: it turns out
    % to be quite tricky to get right.
    %
    \catcode`@=11
    \def\strikethrough#1{\leavevmode\strikethr@ugh#1 \relax}
    \def\strikethr@ugh#1 {\strike@word{#1}\futurelet\next\strike@next}
    \def\strike@next{\ifx\next\relax\else\strike@space\expandafter\strikethr@ugh\fi}
    \def\strike@word#1{\setbox0=\hbox{#1}\rlap{\strike@rule width\wd0}\box0}
    \def\strike@space{\leaders\strike@rule\sp@ce}
    \def\strike@rule{\vrule height2.4pt depth-2pt}
    \def\sp@ce{\hskip\fontdimen2\font plus\fontdimen3\font minus\fontdimen4\font}
    \catcode`@=12
    %
    % (4) Nested lists, enumerated and unenumerated
    % We don't want to have to rely on Plain TeX's `\item` and `\itemitem`.
    % Instead, all items are `\item` (which is redefined as appropriate for the
    % context) and lists begin and end with `\beginenumerate` and `\endenumerate`
    % if they are to be numbered, or `\beginunenumerate` and `\endunenumerate` if not;
    % and those can be fairly arbitrarily nested.
    %
    \catcode`@=11
    \newcount\c@listlevel
    \newcount\c@item
    \newdimen\enumerateindent
    %
    \def\init@str{} % the label of every item
    \def\beginenumerate{%
      \begingroup
      \let\item\enumerate@item
      \advance\c@listlevel\@ne
      \ifnum\c@listlevel=\@ne
        \vskip.5\baselineskip
      \fi
      \ifnum\c@listlevel>\@ne
        \edef\init@str{}%
      \fi
      \c@item\z@
    }
    %
    \def\endenumerate{%
      \par\endgroup
      \ifnum\c@listlevel=\z@
        \vskip.5\baselineskip
      \fi}
    %
    \def\beginunenumerate{%
      \begingroup
      \let\item\unenumerate@item
      \advance\c@listlevel\@ne
      \ifnum\c@listlevel=\@ne
        \vskip.5\baselineskip
      \fi
      \c@item\z@
    }
    %
    \def\endunenumerate{%
      \par\endgroup
      \ifnum\c@listlevel=\z@
        \vskip.5\baselineskip
      \fi}
    %
    \def\enumerate@item{%
      \advance\c@item\@ne
      \enumerateindent=\c@listlevel\parindent
      \ifvmode\else\par\fi
      % if an \item has more than one paragraph only the first one have to be preceded by a label;
      % so \@itemlabel redefine itself the first time is expanded.
      \def\@itemlabel{%
        \llap{\init@str\number\c@item.\quad}%
        \def\@itemlabel{}%
      }%
      \everypar{\setbox0\lastbox
        \hangindent\enumerateindent \hangafter\z@
        \@itemlabel}}
    %
    \def\unenumerate@item{%
      \advance\c@item\@ne
      \enumerateindent=\c@listlevel\parindent
      \ifvmode\else\par\fi
      \everypar{\setbox0\lastbox
        \hangindent\enumerateindent \hangafter\z@
        $\bullet$~}}
    %
    \catcode`@=12
    %
    % (5) Verbatim code samples
    %
    \chardef\other=12
    \def\ttverbatim{\begingroup\ttninepoint\frenchspacing
      \catcode`\\=\other
      \catcode`\{=\other
      \catcode`\}=\other
      \catcode`\$=\other
      \catcode`\&=\other
      \catcode`\#=\other
      \catcode`\%=\other
      \catcode`\~=\other
      \catcode`\_=\other
      \catcode`\^=\other
      \obeyspaces \obeylines \tt}
    %
    \def\|{\leavevmode\hbox{\tt\char`\|}} % vertical line
    \outer\def\begintt{$$\let\par=\endgraf \ttverbatim \parskip=\z@
      \catcode`\|=0 \rightskip-5pc \ttfinish}
    {\catcode`\|=0 |catcode`|\=\other % | is temporary escape character
      |obeylines % end of line is active
      |gdef|ttfinish#1^^M#2\endtt{#1|vbox{#2}|endgroup|tenpoint$$}}
    \catcode`\|=\active
    %
    {\obeylines \gdef|{\ttverbatim \spaceskip\ttglue \let^^M=\  \let|=\endgroup}}
    %
    \catcode`@=11
    \def\beginlines{\par\begingroup\nobreak\medskip\parindent\z@ \obeylines\nobreak \everypar{\strut}}
    \def\beginlinestighter{\par\begingroup\nobreak\parindent=0pt \kern1pt\nobreak \obeylines \everypar{\strut}}
    \def\endlines{\kern1pt\endgroup\medbreak\noindent}
    \catcode`@=12
    %
    % (6) Preform grammar document macros
    % These are no longer really used, but are brief enough to be worth keeping around
    % in case of an underdog comeback story.
    %
    \def\[#1]{\silenttrue\xref|#1|\thinspace{\tt#1}\thinspace} % keyword in syntax
    \def\beginsyntax{\endgraf\nobreak\medskip
      \begingroup \catcode`<=13 \catcode`[=13
      \let\par=\endsyntaxline \obeylines}
    \def\endsyntaxline{\futurelet\next\syntaxswitch}
    \def\syntaxswitch{\ifx\next\<\let\next=\syntaxrule
      \else\ifx\next\endsyntax\let\next=\endgroup
      \else\let\next=\continuerule\fi\fi \next}
    \def\continuerule{\hfil\break\indent\qquad}
    \def\endsyntax{\medbreak\noindent}
    {\catcode`<=13 \catcode`[=13
      \global\let<=\< \global\let[=\[
      \gdef\syntaxrule<#1>{\endgraf\indent\silentfalse\xref\<#1>}}
    \def\is{\ $\longrightarrow$ }
    \def\alt{\ $\vert$ }
    \def\cast#1#2{{#1}$\;\rightarrow\;${#2}}
    \def\comp#1#2{{#1}$\;\sim\;${#2}}
    \def\al{$\alpha$}
    \def\be{$\beta$}
    \def\nonterminal#1{$\langle${\xreffont\pdfliteral direct{1 1 0 0 k}#1\special{PDF:0 g}}$\rangle$}
    \def\fixed#1{{\bf #1}}
    \def\from{\par\qquad\quad$\leftarrow$\quad}
    \def\continuation{\par\qquad\quad\phantom{$\leftarrow$}\quad}
    %
    % (7) Miscellaneous math macros which seem to be expected by MathJax users,
    % even though they are not in Plain TeX
    %
    \def\frac#1#2{{{#1}\over{#2}}}
    \def\cfrac#1#2{{{#1}\over{#2}}}
    \def\dfrac#1#2{{{#1}\over{#2}}}
    \def\tfrac#1#2{{{#1}\over{#2}}}
    \def\mathit#1{{\hbox{\it #1}}}
    \def\mathrm#1{{\hbox{\rm #1}}}
    \def\mathbf#1{{\hbox{\bf #1}}}
    \def\text#1{{\hbox{#1}}}
    \def\lvert{|}
    \def\rvert{|}
    \def\implies{\Rightarrow}
    %
    % (8) And finally, something other than macros and fonts: the page design.
    % First, header and footer:
    %
    \nopagenumbers
    \newif\iftitle
    \def\titlepage{\global\titletrue}
    \headline={\iftitle\hfil\global\titlefalse\else\tenrm\hfil{\sourcefont \firstmark\qquad\the\pageno}\fi}
    \titlepage
    %
    % Tolerances for spacing:
    %
    \raggedbottom
    \tolerance=20000
    %
    % End of `inweb-macros.tex`


% Weave of 'Complete Program' generated by inweb
\weavesectionsss{}{Chapter 1\quad Layout}{}{\P}{1/ld}%
\medskip
\smallskip\noindent {\bf Layout Demonstration}\qquad
This section exercises paragraph structure, direct notations for inserted gadgets, and footnotes.\weavesectionss{}{Layout Demonstration}{}{\P}{}%
\smallskip\par\noindent{\it This section exercises paragraph structure, direct notations for inserted gadgets, and footnotes.}\smallskip\noindent
\medskip\hrule\smallskip\par\noindent\usagefont {\inweblinktoanchor{para1_ld1}{$\S$1}}~Disclaimer; {\inweblinktoanchor{para1_ld2}{$\S$2}}~Named paragraph; {\inweblinktoanchor{para1_ld6}{$\S$6}}~Another named paragraph; {\inweblinktoanchor{para1_ld7}{$\S$7}}~Gadgetry}\par\medskip\hrule\bigskip
\weavesections{1}{Disclaimer}{Chapter 1: Layout - Layout Demonstration\quad$\S$1}{\S}{1/ld}%
\inwebanchor{para1_ld1}The content of this web is essentially meaningless. It exists only to test
the renderers.


\weavesections{2}{Named paragraph}{Chapter 1: Layout - Layout Demonstration\quad$\S$2}{\S}{1/ld}%
\inwebanchor{para1_ld2}This paragraph has a name, and is usually rendered with a modest subheading.\footnote{${}^{1}$}{Also two footnotes.

}
There is very little else to it.\footnote{${}^{2}$}{Save for the formula $n\mapsto 3n + 1$.

}


\weavesection{3}{}{Chapter 1: Layout - Layout Demonstration\quad$\S$3}{\S}{1/ld}%
\inwebanchor{para1_ld3}This paragraph is anonymous, and will contain some definitions.


\beginlinestighter
\hskip1em \noindent \constantsyntaxcolouring{}|ONE|\plainsyntaxcolouring{} $\equiv$ ||\constantsyntaxcolouring{}|1|\endcolouring{}||
\hskip1em \noindent \constantsyntaxcolouring{}|MULTIPLE|\plainsyntaxcolouring{} $\equiv$ ||\constantsyntaxcolouring{}|3|\endcolouring{}||
\hskip1em \noindent \constantsyntaxcolouring{}|MULTIPLE|\plainsyntaxcolouring{} $\equiv_{\hbox{\defnannotationfont def}}$ ||\constantsyntaxcolouring{}|10|\endcolouring{}||
\hskip1em \noindent {\defnannotationfont enum} \constantsyntaxcolouring{}|RED_COLOUR|\plainsyntaxcolouring{}  ||\constantsyntaxcolouring{}|1|\endcolouring{}||
\hskip1em \noindent {\defnannotationfont enum} \constantsyntaxcolouring{}|GREEN_COLOUR|\plainsyntaxcolouring{}  
\hskip1em \noindent {\defnannotationfont enum} \constantsyntaxcolouring{}|BLUE_COLOUR|\plainsyntaxcolouring{}  
\endlines
\weavesection{4}{}{Chapter 1: Layout - Layout Demonstration\quad$\S$4}{\S}{1/ld}%
\inwebanchor{para1_ld4}The trivial function at {\inweblinktoanchor{s1line12}{A}} implements the famous iteration of
Lothar Collatz (1910-1990). The use of |MULTIPLE| at {\inweblinktoanchor{s1line4}{B}} triples and adds 1.


\beginlines
\hskip1em \inwebanchor{s1line1}|A |{$\Rightarrow$}\hskip1em ||\identifiersyntaxcolouring{}|int|\plainsyntaxcolouring{}| |\endcolouring{}||||\functionsyntaxcolouring{}|collatz|\plainsyntaxcolouring{}||||\plainsyntaxcolouring{}|(|\identifiersyntaxcolouring{}|int|\plainsyntaxcolouring{}| |\identifiersyntaxcolouring{}|x|\plainsyntaxcolouring{}|) {            |\endcolouring{}||||\commentsyntaxcolouring{}|/* A */|\plainsyntaxcolouring{}||
\hskip1em |  |\phantom{$\Rightarrow$}\hskip1em ||\plainsyntaxcolouring{}|    |\endcolouring{}||\inweblinktopara{para105}{\definitionsyntaxcolouring{}$\langle$\holonnametext{Print out x} \holonnumbertext{4.1}\endcolouring{}$\rangle$ }
\hskip1em |  |\phantom{$\Rightarrow$}\hskip1em ||\plainsyntaxcolouring{}|    |\reservedsyntaxcolouring{}|if|\plainsyntaxcolouring{}| (|\identifiersyntaxcolouring{}|x|\plainsyntaxcolouring{}| % |\constantsyntaxcolouring{}|2|\plainsyntaxcolouring{}| == |\constantsyntaxcolouring{}|0|\plainsyntaxcolouring{}|) |\reservedsyntaxcolouring{}|return|\plainsyntaxcolouring{}| |\identifiersyntaxcolouring{}|x|\plainsyntaxcolouring{}|/2;|\endcolouring{}||
\hskip1em \inwebanchor{s1line4}|B |{$\Rightarrow$}\hskip1em ||\plainsyntaxcolouring{}|    |\reservedsyntaxcolouring{}|return|\plainsyntaxcolouring{}| |\constantsyntaxcolouring{}|MULTIPLE|\plainsyntaxcolouring{}|*|\identifiersyntaxcolouring{}|x|\plainsyntaxcolouring{}| + |\constantsyntaxcolouring{}|1|\plainsyntaxcolouring{}|;      |\endcolouring{}||||\commentsyntaxcolouring{}|/* B */|\plainsyntaxcolouring{}||
\hskip1em |  |\phantom{$\Rightarrow$}\hskip1em ||\plainsyntaxcolouring{}|}|\endcolouring{}||
\endlines
\weavesection{4.1}{}{Chapter 1: Layout - Layout Demonstration\quad$\S$4.1}{\S}{1/ld}%
\inwebanchor{para1_ld4.1}A named holon appears next. This one appears after some commentary, but
that is not obligatory, as the next will show.


\par\smallskip\noindent
\inwebanchor{para105}\definitionsyntaxcolouring{}$\langle$\holonnametext{Print out x} \holonnumbertext{4.1}\endcolouring{}$\rangle$ $\equiv$
\beginlinestighter
\hskip1em ||\plainsyntaxcolouring{}|    |\identifiersyntaxcolouring{}|printf|\plainsyntaxcolouring{}|(|\stringsyntaxcolouring{}|"x = %d\n"|\plainsyntaxcolouring{}|, |\identifiersyntaxcolouring{}|x|\plainsyntaxcolouring{}|);   |\endcolouring{}||||\commentsyntaxcolouring{}|/* `x` here is what we called $n$ above */|\plainsyntaxcolouring{}||
\hskip1em ||\plainsyntaxcolouring{}|    |\endcolouring{}||||\charactersyntaxcolouring{}|sleep(1);|\endcolouring{}||||\plainsyntaxcolouring{}|                |\endcolouring{}||||\commentsyntaxcolouring{}|// that was done with a verbatim splice|\plainsyntaxcolouring{}||
\endlines
\par\noindent\penalty10000
\endnotetext{This code is continued in {\inweblinktoanchor{para1_ld5}{$\S$5}}, and used in {\inweblinktoanchor{para1_ld4}{$\S$4}}.}\smallskip
\weavesection{5}{}{Chapter 1: Layout - Layout Demonstration\quad$\S$5}{\S}{1/ld}%
\inwebanchor{para1_ld5}\inweblinktopara{para105}{\definitionsyntaxcolouring{}$\langle$\holonnametext{Print out x} \holonnumbertext{4.1}\endcolouring{}$\rangle$ $+${}$\equiv$}
\beginlinestighter
\hskip1em ||\plainsyntaxcolouring{}|    |\identifiersyntaxcolouring{}|print|\plainsyntaxcolouring{}|(|\stringsyntaxcolouring{}|"x is still %d\n"|\plainsyntaxcolouring{}|, |\identifiersyntaxcolouring{}|x|\plainsyntaxcolouring{}|);|\endcolouring{}||
\endlines
\par\noindent\penalty10000
\endnotetext{This code is a continuation of {\inweblinktoanchor{para1_ld4.1}{$\S$4.1}}.}\smallskip
\weavesections{6}{Another named paragraph}{Chapter 1: Layout - Layout Demonstration\quad$\S$6}{\S}{1/ld}%
\inwebanchor{para1_ld6}Note the duplicate label here: the link {\inweblinktoanchor{s1line12}{A}} goes to the one in this function.


\beginlines
\hskip1em \inwebanchor{s1line12}|A |{$\Rightarrow$}\hskip1em ||\identifiersyntaxcolouring{}|int|\plainsyntaxcolouring{}| |\endcolouring{}||||\functionsyntaxcolouring{}|lothar|\plainsyntaxcolouring{}||||\plainsyntaxcolouring{}|(|\identifiersyntaxcolouring{}|int|\plainsyntaxcolouring{}| |\identifiersyntaxcolouring{}|n|\plainsyntaxcolouring{}|) { |\endcolouring{}||||\commentsyntaxcolouring{}|/* A */|\plainsyntaxcolouring{}||
\hskip1em |  |\phantom{$\Rightarrow$}\hskip1em ||\plainsyntaxcolouring{}|    |\reservedsyntaxcolouring{}|return|\plainsyntaxcolouring{}| |\endcolouring{}||||\functionsyntaxcolouring{}|collatz|\endcolouring{}||||\plainsyntaxcolouring{}|(|\identifiersyntaxcolouring{}|n|\plainsyntaxcolouring{}|);|\endcolouring{}||
\hskip1em |  |\phantom{$\Rightarrow$}\hskip1em ||\plainsyntaxcolouring{}|}|\endcolouring{}||
\endlines
\weavesections{7}{Gadgetry}{Chapter 1: Layout - Layout Demonstration\quad$\S$7}{\S}{1/ld}%
\inwebanchor{para1_ld7}This uses notation for insertions by hand, so to speak, rather than via a
Markdown extension, so the following involves non-commentary chunks in the
literate source tree.

\inwebimagewidth{inweb/Tests/omnium/Figures/stlucia.jpg}{15}

\weavesection{8}{}{Chapter 1: Layout - Layout Demonstration\quad$\S$8}{\S}{1/ld}%
\inwebanchor{para1_ld8}If this renders to a format which supports audio, then this will be the call
of Steller's Jay, from the Cornell Lab of Ornithology collection
{\inweblinktourl{https://dl.allaboutbirds.org}{Voices of Western Backyard Birds}}:


\weavesection{9}{}{Chapter 1: Layout - Layout Demonstration\quad$\S$9}{\S}{1/ld}%
\inwebanchor{para1_ld9}This doesn't look much (and isn't): a raw piece of HTML pasted in from a
file, and visible only when this web is rendered to HTML.


\weavesection{10}{}{Chapter 1: Layout - Layout Demonstration\quad$\S$10}{\S}{1/ld}%
\inwebanchor{para1_ld10}A modest slide carousel, again using direct insertions:



\medskip\hrule\smallskip
\beginlines
\hskip1em||\plainsyntaxcolouring{}|    |\functionsyntaxcolouring{}|ROOT|\plainsyntaxcolouring{}| ---> |\functionsyntaxcolouring{}|DOCUMENT|\plainsyntaxcolouring{}||
\hskip1em||\endcolouring{}||\endlines

\centerline{\bf 1/3. Stage 1 - Raw tree}


\smallskip\hrule\medskip

\medskip\hrule\smallskip
\centerline{\bf 2/3. Stage 2 - Developed tree}

\beginlines
\hskip1em||\plainsyntaxcolouring{}|    |\functionsyntaxcolouring{}|ROOT|\plainsyntaxcolouring{}| ---> |\functionsyntaxcolouring{}|DOCUMENT|\plainsyntaxcolouring{}||
\hskip1em|                |{\tt\|}||
\hskip1em|              |\functionsyntaxcolouring{}|NODE|\plainsyntaxcolouring{}| |\constantsyntaxcolouring{}|1|\plainsyntaxcolouring{}|  ---  |\functionsyntaxcolouring{}|NODE|\plainsyntaxcolouring{}| |\constantsyntaxcolouring{}|2|\plainsyntaxcolouring{}|  ---  |\functionsyntaxcolouring{}|NODE|\plainsyntaxcolouring{}| |\constantsyntaxcolouring{}|3|\plainsyntaxcolouring{}|  --- ...|
\hskip1em||\endcolouring{}||\endlines


\smallskip\hrule\medskip

\medskip\hrule\smallskip
\beginlines
\hskip1em||\plainsyntaxcolouring{}|    |\functionsyntaxcolouring{}|ROOT|\plainsyntaxcolouring{}| ---> |\functionsyntaxcolouring{}|DOCUMENT|\plainsyntaxcolouring{}||
\hskip1em|                |{\tt\|}||
\hskip1em|              |\functionsyntaxcolouring{}|NODE|\plainsyntaxcolouring{}| |\constantsyntaxcolouring{}|1|\plainsyntaxcolouring{}|  ---  |\functionsyntaxcolouring{}|NODE|\plainsyntaxcolouring{}| |\constantsyntaxcolouring{}|2|\plainsyntaxcolouring{}|  ---  |\functionsyntaxcolouring{}|NODE|\plainsyntaxcolouring{}| |\constantsyntaxcolouring{}|3|\plainsyntaxcolouring{}|  --- ...|
\hskip1em|                |{\tt\|}|            |{\tt\|}|            |{\tt\|}||
\hskip1em|              |\elementsyntaxcolouring{}|text|\plainsyntaxcolouring{}| |\constantsyntaxcolouring{}|1|\plainsyntaxcolouring{}|       |\elementsyntaxcolouring{}|text|\plainsyntaxcolouring{}| |\constantsyntaxcolouring{}|2|\plainsyntaxcolouring{}|       |\elementsyntaxcolouring{}|text|\plainsyntaxcolouring{}| |\constantsyntaxcolouring{}|3|\plainsyntaxcolouring{}|  ...|
\hskip1em||\endcolouring{}||\endlines

\centerline{\bf 3/3. Stage 3 - Completed tree}


\smallskip\hrule\medskip
\weavesectionsss{}{Chapter 2\quad Commentary}{}{\P}{2/md}%
\medskip
\smallskip\noindent {\bf Markdown Demonstration}\qquad
This section exercises Markdown-format commentary, complete with its many Inweb extensions. These provide an alternative way to insert gadgets.\weavesectionss{}{Markdown Demonstration}{}{\P}{}%
\smallskip\par\noindent{\it This section exercises Markdown-format commentary, complete with its many Inweb extensions. These provide an alternative way to insert gadgets.}\smallskip\noindent
\medskip\hrule\smallskip\par\noindent\usagefont {\inweblinktoanchor{para2_md1}{$\S$1}}~Commentary; {\inweblinktoanchor{para2_md2}{$\S$2}}~Character set; {\inweblinktoanchor{para2_md3}{$\S$3}}~ISO Latin-1; {\inweblinktoanchor{para2_md4}{$\S$4}}~Mathematics; {\inweblinktoanchor{para2_md5}{$\S$5}}~Internal subheadings; {\inweblinktoanchor{para2_md6}{$\S$6}}~Lists and tables; {\inweblinktoanchor{para2_md7}{$\S$7}}~Block quotations; {\inweblinktoanchor{para2_md8}{$\S$8}}~Links; {\inweblinktoanchor{para2_md9}{$\S$9}}~Code extracts; {\inweblinktoanchor{para2_md10}{$\S$10}}~Images and gadgets; {\inweblinktoanchor{para2_md11}{$\S$11}}~Carousels; {\inweblinktoanchor{para2_md12}{$\S$12}}~A gratuitously long table; {\inweblinktoanchor{para2_md13}{$\S$13}}~Index}\par\medskip\hrule\bigskip
\weavesections{1}{Commentary}{Chapter 2: Commentary - Markdown Demonstration\quad$\S$1}{\S}{2/md}%
\inwebanchor{para2_md1}The content of this section is again meaningless. It exists only to test
the renderers.

Here is some {\it emphasis} and {\bf strong emphasis} and \strikethrough{text struck through}.

Here is some |literal backticked text|.

Some index entries appear here, invisibly.


\weavesections{2}{Character set}{Chapter 2: Commentary - Markdown Demonstration\quad$\S$2}{\S}{2/md}%
\inwebanchor{para2_md2}These are potentially tricky for TeX: $\{$Knuth$\}$. A$\backslash$B. A\_B. A\#B. \$40. A\|B. 30\%. 2\char`^56.

And in backticked literal text: |{Knuth}. A\B. A_B. A#B. $40. A|{\tt\|}|B. 30%. 2^26.|

These are potentially tricky for HTML: {\tt<}, \&, {\tt>}.

And in backticked literal text: |<, &, >|.

``Double quoted.'' `Single quoted.'

In TeX variants, it's `reasonable' to hope for ``smart quotes'' to be imposed, but
more questionable whether phrases like église château mañana will be accented.


\weavesections{3}{ISO Latin-1}{Chapter 2: Commentary - Markdown Demonstration\quad$\S$3}{\S}{2/md}%
\inwebanchor{para2_md3}Some formats can take most of Unicode, but others are more limited, and plain TeX
is stuck with a small selection of ASCII, leaving the following ragged. It's in
both tabular and untabular form because of issues with TeX rendering of table cells,
which mean that implementation of magic characters works differently in those cases.


\medskip\centerline{\vbox{\offinterlineskip
\hrule
\halign{&\vrule#&\strut\quad\hfil#\hfil\quad&\vrule#&\strut\quad\hfil#\hfil\quad&\vrule#&\strut\quad\hfil#\hfil\quad&\vrule#&\strut\quad\hfil#\hfil\quad&\vrule#&\strut\quad\hfil#\hfil\quad&\vrule#&\strut\quad\hfil#\hfil\quad&\vrule#&\strut\quad\hfil#\hfil\quad&\vrule#&\strut\quad\hfil#\hfil\quad&\vrule#&\strut\quad\hfil#\hfil\quad&\vrule#&\strut\quad\hfil#\hfil\quad&\vrule#&\strut\quad\hfil#\hfil\quad&\vrule#&\strut\quad\hfil#\hfil\quad&\vrule#&\strut\quad\hfil#\hfil\quad&\vrule#&\strut\quad\hfil#\hfil\quad&\vrule#&\strut\quad\hfil#\hfil\quad&\vrule#&\strut\quad\hfil#\hfil\quad&\vrule#&\strut\quad\hfil#\hfil\quad\cr
height2pt&\omit&&\omit&&\omit&&\omit&&\omit&&\omit&&\omit&&\omit&&\omit&&\omit&&\omit&&\omit&&\omit&&\omit&&\omit&&\omit&&\omit&\cr
&&&0&&1&&2&&3&&4&&5&&6&&7&&8&&9&&A&&B&&C&&D&&E&&F&\cr
height2pt&\omit&&\omit&&\omit&&\omit&&\omit&&\omit&&\omit&&\omit&&\omit&&\omit&&\omit&&\omit&&\omit&&\omit&&\omit&&\omit&&\omit&\cr
\noalign{\hrule}
height2pt&\omit&&\omit&&\omit&&\omit&&\omit&&\omit&&\omit&&\omit&&\omit&&\omit&&\omit&&\omit&&\omit&&\omit&&\omit&&\omit&&\omit&\cr
&2x&&SP&&!&&``&&\#&&\$&&\%&&\&&&`&&(&&)&&*&&+&&,&&-&&.&&/&\cr
&3x&&0&&1&&2&&3&&4&&5&&6&&7&&8&&9&&:&&;&&{\tt<}&&=&&{\tt>}&&?&\cr
&4x&&@&&A&&B&&C&&D&&E&&F&&G&&H&&I&&J&&K&&L&&M&&N&&O&\cr
&5x&&P&&Q&&R&&S&&T&&U&&V&&W&&X&&Y&&Z&&[&&$\backslash$&&]&&\char`^&&\_&\cr
&6x&&`&&a&&b&&c&&d&&e&&f&&g&&h&&i&&j&&k&&l&&m&&n&&o&\cr
&7x&&p&&q&&r&&s&&t&&u&&v&&w&&x&&y&&z&&$\{$&&\|&&$\}$&&\char`~&&&\cr
&Ax&&NB&&¡&&¢&&£&&¤&&¥&&¦&&§&&¨&&©&&ª&&«&&¬&&SH&&®&&¯&\cr
&Bx&&°&&±&&²&&³&&´&&µ&&¶&&·&&¸&&¹&&º&&»&&¼&&½&&¾&&¿&\cr
&Cx&&À&&Á&&Â&&Ã&&Ä&&Å&&Æ&&Ç&&È&&É&&Ê&&Ë&&Ì&&Í&&Î&&Ï&\cr
&Dx&&Ð&&Ñ&&Ò&&Ó&&Ô&&Õ&&Ö&&×&&Ø&&Ù&&Ú&&Û&&Ü&&Ý&&Þ&&ß&\cr
&Ex&&à&&á&&â&&ã&&ä&&å&&æ&&ç&&è&&é&&ê&&ë&&ì&&í&&î&&ï&\cr
&Fx&&ð&&ñ&&ò&&ó&&ô&&õ&&ö&&÷&&ø&&ù&&ú&&û&&ü&&ý&&þ&&ÿ&\cr
height2pt&\omit&&\omit&&\omit&&\omit&&\omit&&\omit&&\omit&&\omit&&\omit&&\omit&&\omit&&\omit&&\omit&&\omit&&\omit&&\omit&&\omit&\cr
}\hrule}}\medskip


\medskip\centerline{\vbox{\offinterlineskip
\hrule
\halign{&\vrule#&\strut\quad\hfil#\hfil\quad&\vrule#&\strut\quad\hfil#\hfil\quad&\vrule#&\strut\quad\hfil#\hfil\quad&\vrule#&\strut\quad\hfil#\hfil\quad&\vrule#&\strut\quad\hfil#\hfil\quad&\vrule#&\strut\quad\hfil#\hfil\quad&\vrule#&\strut\quad\hfil#\hfil\quad&\vrule#&\strut\quad\hfil#\hfil\quad&\vrule#&\strut\quad\hfil#\hfil\quad&\vrule#&\strut\quad\hfil#\hfil\quad&\vrule#&\strut\quad\hfil#\hfil\quad&\vrule#&\strut\quad\hfil#\hfil\quad&\vrule#&\strut\quad\hfil#\hfil\quad&\vrule#&\strut\quad\hfil#\hfil\quad&\vrule#&\strut\quad\hfil#\hfil\quad&\vrule#&\strut\quad\hfil#\hfil\quad&\vrule#&\strut\quad\hfil#\hfil\quad\cr
height2pt&\omit&&\omit&&\omit&&\omit&&\omit&&\omit&&\omit&&\omit&&\omit&&\omit&&\omit&&\omit&&\omit&&\omit&&\omit&&\omit&&\omit&\cr
&&&0&&1&&2&&3&&4&&5&&6&&7&&8&&9&&A&&B&&C&&D&&E&&F&\cr
height2pt&\omit&&\omit&&\omit&&\omit&&\omit&&\omit&&\omit&&\omit&&\omit&&\omit&&\omit&&\omit&&\omit&&\omit&&\omit&&\omit&&\omit&\cr
\noalign{\hrule}
height2pt&\omit&&\omit&&\omit&&\omit&&\omit&&\omit&&\omit&&\omit&&\omit&&\omit&&\omit&&\omit&&\omit&&\omit&&\omit&&\omit&&\omit&\cr
&2x&&|SP|&&|!|&&|“|&&|\#|&&|\$|&&|\%|&&|\&|&&|‘|&&|(|&&|)|&&|*|&&|+|&&|,|&&|-|&&|.|&&|/|&\cr
&3x&&|0|&&|1|&&|2|&&|3|&&|4|&&|5|&&|6|&&|7|&&|8|&&|9|&&|:|&&|;|&&|<|&&|=|&&|>|&&|?|&\cr
&4x&&|@|&&|A|&&|B|&&|C|&&|D|&&|E|&&|F|&&|G|&&|H|&&|I|&&|J|&&|K|&&|L|&&|M|&&|N|&&|O|&\cr
&5x&&|P|&&|Q|&&|R|&&|S|&&|T|&&|U|&&|V|&&|W|&&|X|&&|Y|&&|Z|&&|[|&&||$\backslash$||&&|]|&&|\char`^|&&|\_|&\cr
&6x&&|`|&&|a|&&|b|&&|c|&&|d|&&|e|&&|f|&&|g|&&|h|&&|i|&&|j|&&|k|&&|l|&&|m|&&|n|&&|o|&\cr
&7x&&|p|&&|q|&&|r|&&|s|&&|t|&&|u|&&|v|&&|w|&&|x|&&|y|&&|z|&&||$\{$||&&bar&&||$\}$||&&||\tt\char`~||&&&\cr
&Ax&&|NB|&&|¡|&&|¢|&&|£|&&|¤|&&|¥|&&|¦|&&|§|&&|¨|&&|©|&&|ª|&&|«|&&|¬|&&|SH|&&|®|&&|¯|&\cr
&Bx&&|°|&&|±|&&|²|&&|³|&&|´|&&|µ|&&|¶|&&|·|&&|¸|&&|¹|&&|º|&&|»|&&|¼|&&|½|&&|¾|&&|¿|&\cr
&Cx&&|À|&&|Á|&&|Â|&&|Ã|&&|Ä|&&|Å|&&|Æ|&&|Ç|&&|È|&&|É|&&|Ê|&&|Ë|&&|Ì|&&|Í|&&|Î|&&|Ï|&\cr
&Dx&&|Ð|&&|Ñ|&&|Ò|&&|Ó|&&|Ô|&&|Õ|&&|Ö|&&|×|&&|Ø|&&|Ù|&&|Ú|&&|Û|&&|Ü|&&|Ý|&&|Þ|&&|ß|&\cr
&Ex&&|à|&&|á|&&|â|&&|ã|&&|ä|&&|å|&&|æ|&&|ç|&&|è|&&|é|&&|ê|&&|ë|&&|ì|&&|í|&&|î|&&|ï|&\cr
&Fx&&|ð|&&|ñ|&&|ò|&&|ó|&&|ô|&&|õ|&&|ö|&&|÷|&&|ø|&&|ù|&&|ú|&&|û|&&|ü|&&|ý|&&|þ|&&|ÿ|&\cr
height2pt&\omit&&\omit&&\omit&&\omit&&\omit&&\omit&&\omit&&\omit&&\omit&&\omit&&\omit&&\omit&&\omit&&\omit&&\omit&&\omit&&\omit&\cr
}\hrule}}\medskip

SP   !   ``   \#   \$   \%   \&   `   (   )   *   +   ,   -   .   /

0   1   2   3   4   5   6   7   8   9   :   ;   {\tt<}   =   {\tt>}   ?

@   A   B   C   D   E   F   G   H   I   J   K   L   M   N   O

P   Q   R   S   T   U   V   W   X   Y   Z   [   $\backslash$   ]   \char`^   \_

`   a   b   c   d   e   f   g   h   i   j   k   l   m   n   o

p   q   r   s   t   u   v   w   x   y   z   $\{$   \|   $\}$   \char`~

NB   ¡   ¢   £   ¤   ¥   ¦   §   ¨   ©   ª   «   ¬  SH   ®   ¯

°   ±   ²   ³   ´   µ   ¶   ·   ¸   ¹   º   »   ¼   ½   ¾   ¿

À   Á   Â   Ã   Ä   Å   Æ   Ç   È   É   Ê   Ë   Ì   Í   Î   Ï

Ð   Ñ   Ò   Ó   Ô   Õ   Ö   ×   Ø   Ù   Ú   Û   Ü   Ý   Þ   ß

à   á   â   ã   ä   å   æ   ç   è   é   ê   ë   ì   í   î   ï

ð   ñ   ò   ó   ô   õ   ö   ÷   ø   ù   ú   û   ü   ý   þ   ÿ

|SP|   |!|   |“|   |#|   |$|   |%|   |&|   |‘|   |(|   |)|   |*|   |+|   |,|   |-|   |.|   |/|

|0|   |1|   |2|   |3|   |4|   |5|   |6|   |7|   |8|   |9|   |:|   |;|   |<|   |=|   |>|   |?|

|@|   |A|   |B|   |C|   |D|   |E|   |F|   |G|   |H|   |I|   |J|   |K|   |L|   |M|   |N|   |O|

|P|   |Q|   |R|   |S|   |T|   |U|   |V|   |W|   |X|   |Y|   |Z|   |[|   |\|   |]|   |^|   |_|

|`|   |a|   |b|   |c|   |d|   |e|   |f|   |g|   |h|   |i|   |j|   |k|   |l|   |m|   |n|   |o|

|p|   |q|   |r|   |s|   |t|   |u|   |v|   |w|   |x|   |y|   |z|   |{|   ||{\tt\|}||   |}|   |~|

|NB|   |¡|   |¢|   |£|   |¤|   |¥|   |¦|   |§|   |¨|   |©|   |ª|   |«|   |¬|  |SH|   |®|   |¯|

|°|   |±|   |²|   |³|   |´|   |µ|   |¶|   |·|   |¸|   |¹|   |º|   |»|   |¼|   |½|   |¾|   |¿|

|À|   |Á|   |Â|   |Ã|   |Ä|   |Å|   |Æ|   |Ç|   |È|   |É|   |Ê|   |Ë|   |Ì|   |Í|   |Î|   |Ï|

|Ð|   |Ñ|   |Ò|   |Ó|   |Ô|   |Õ|   |Ö|   |×|   |Ø|   |Ù|   |Ú|   |Û|   |Ü|   |Ý|   |Þ|   |ß|

|à|   |á|   |â|   |ã|   |ä|   |å|   |æ|   |ç|   |è|   |é|   |ê|   |ë|   |ì|   |í|   |î|   |ï|

|ð|   |ñ|   |ò|   |ó|   |ô|   |õ|   |ö|   |÷|   |ø|   |ù|   |ú|   |û|   |ü|   |ý|   |þ|   |ÿ|


\weavesections{4}{Mathematics}{Chapter 2: Commentary - Markdown Demonstration\quad$\S$4}{\S}{2/md}%
\inwebanchor{para2_md4}Zolotarev's Lemma says that $a\neq 0$ has a square root modulo a prime $p$
if and only if the rearrangement of $\lbrace 0, 1, 2, \dots, p-1\rbrace$ produced by multiplying
by $a$ is an even permutation. That is, if $\sigma(x) = ax {\rm ~mod~} p$, then
$$ \left({\frac {a}{p}}\right) = \cases{1 & if $\sigma$ even,\cr
-1 & if $\sigma$ odd.\cr} $$


\weavesections{5}{Internal subheadings}{Chapter 2: Commentary - Markdown Demonstration\quad$\S$5}{\S}{2/md}%
\inwebanchor{para2_md5}
\medskip\par\noindent{\bf Subhead level 1}


\medskip\par\noindent{\bf Subhead level 2}


\medskip\par\noindent{\it Subhead level 3}

This seems a good moment for what Markdown calls a thematic break.


\bigskip\hrule\bigskip

And here are some more invisible entries.


\weavesections{6}{Lists and tables}{Chapter 2: Commentary - Markdown Demonstration\quad$\S$6}{\S}{2/md}%
\inwebanchor{para2_md6}Some notable etymologies:

\beginunenumerate

\item \smallskip
Lanthanum from the Greek ``lanthanein'', meaning to be hidden.



\item \smallskip
Praseodymium from the Greek ``prasios'', meaning leek-green, and ``didymos'', meaning twin.



\item \smallskip
All of these from the mining village of Ytterby, on the Swedish island of Resaro:

\beginenumerate

\item Terbium

\item Erbium

\item Yttrium

\item Ytterbium
\endenumerate

\endunenumerate
To tabulate:


\medskip\centerline{\vbox{\offinterlineskip
\hrule
\halign{&\vrule#&\strut\quad#\hfil\quad&\vrule#&\strut\quad\hfil#\hfil\quad&\vrule#&\strut\quad\hfil#\quad\cr
height2pt&\omit&&\omit&&\omit&\cr
&Element&&Name&&Usages&\cr
height2pt&\omit&&\omit&&\omit&\cr
\noalign{\hrule}
height2pt&\omit&&\omit&&\omit&\cr
&21 Sc&&Scandium&&aerospace alloys&\cr
&39 Y&&Yttrium&&tooth crowns, jet engines, light bulbs&\cr
&57 La&&Lanthanum&&glass additive, lenses&\cr
&58 Ce&&Cerium&&polishing powder, self-cleaning ovens&\cr
&59 Pr&&Praseodymium&&lasers, arc lights&\cr
height2pt&\omit&&\omit&&\omit&\cr
}\hrule}}\medskip


\weavesections{7}{Block quotations}{Chapter 2: Commentary - Markdown Demonstration\quad$\S$7}{\S}{2/md}%
\inwebanchor{para2_md7}\smallskip\par{\narrower\narrower\noindent
In 1787, Lieutenant Carl Axel Arrhenius found an unidentified black mineral
in the Ytterby mine, whose purpose was to produce quartz and later feldspar.

\smallskip}
Later on:

\smallskip\par{\narrower\narrower\noindent
In 1953, the mine was renovated and used for the storage of jet fuel, MC 77,
but the storage method led to contamination.

\smallskip}
We also support alerts:

{\narrower\smallskip\noindent
$\otimes$ {\bf Warning}

\noindent Passenger ships of the {\it Waxholmsbolaget} do not call at Ytterby pier in the
winter months.

\smallskip}

\weavesections{8}{Links}{Chapter 2: Commentary - Markdown Demonstration\quad$\S$8}{\S}{2/md}%
\inwebanchor{para2_md8}Regular Markdown links first. Here is a {\inweblinktourl{https://en.wikipedia.org/wiki/Acorn_System_BASIC}{typical link without problematic
characters in its URI}}, and a
{\inweblinktourl{https://en.wikipedia.org/wiki/D\%C3\%A9class\%C3\%A9e}{link with an accent in}}, and lastly
a {\inweblinktourl{https://en.wikipedia.org/w/index.php?title=\%E2\%84\%B5}{third link with a Hebrew letter}}.

Inweb has its own slashed links, which we try next.
First {\inweblinktourl{https://en.wikipedia.org}{an external slashed link to Wikipedia}},
and then {\inweblinktoanchor{s1line12}{an internal one to code label A}}.


\weavesections{9}{Code extracts}{Chapter 2: Commentary - Markdown Demonstration\quad$\S$9}{\S}{2/md}%
\inwebanchor{para2_md9}To begin with, an unfenced code extract:

\beginlines
\hskip1em||\identifiersyntaxcolouring{}|factorial|\plainsyntaxcolouring{}| = |\constantsyntaxcolouring{}|1|\plainsyntaxcolouring{}||
\hskip1em||\reservedsyntaxcolouring{}|for|\plainsyntaxcolouring{}| |\identifiersyntaxcolouring{}|i|\plainsyntaxcolouring{}| |\identifiersyntaxcolouring{}|in|\plainsyntaxcolouring{}| |\identifiersyntaxcolouring{}|range|\plainsyntaxcolouring{}|(2, |\identifiersyntaxcolouring{}|n|\plainsyntaxcolouring{}| + |\constantsyntaxcolouring{}|1|\plainsyntaxcolouring{}|):|
\hskip1em|    |\identifiersyntaxcolouring{}|factorial|\plainsyntaxcolouring{}| *= |\identifiersyntaxcolouring{}|i|\plainsyntaxcolouring{}||
\hskip1em||
\hskip1em||\identifiersyntaxcolouring{}|print|\plainsyntaxcolouring{}|(|\identifiersyntaxcolouring{}|factorial|\plainsyntaxcolouring{}|)|
\hskip1em||\endcolouring{}||\endlines
And the same extract fenced:

\beginlines
\hskip1em||\identifiersyntaxcolouring{}|factorial|\plainsyntaxcolouring{}| = |\constantsyntaxcolouring{}|1|\plainsyntaxcolouring{}||
\hskip1em||\reservedsyntaxcolouring{}|for|\plainsyntaxcolouring{}| |\identifiersyntaxcolouring{}|i|\plainsyntaxcolouring{}| |\identifiersyntaxcolouring{}|in|\plainsyntaxcolouring{}| |\identifiersyntaxcolouring{}|range|\plainsyntaxcolouring{}|(2, |\identifiersyntaxcolouring{}|n|\plainsyntaxcolouring{}| + |\constantsyntaxcolouring{}|1|\plainsyntaxcolouring{}|):|
\hskip1em|    |\identifiersyntaxcolouring{}|factorial|\plainsyntaxcolouring{}| *= |\identifiersyntaxcolouring{}|i|\plainsyntaxcolouring{}||
\hskip1em||
\hskip1em||\identifiersyntaxcolouring{}|print|\plainsyntaxcolouring{}|(|\identifiersyntaxcolouring{}|factorial|\plainsyntaxcolouring{}|)|
\hskip1em||\endcolouring{}||\endlines
This time syntax-coloured as None:

\beginlines
\hskip1em||\plainsyntaxcolouring{}|factorial = 1|
\hskip1em|for i in range(2, n + 1):|
\hskip1em|    factorial *= i|
\hskip1em||
\hskip1em|print(factorial)|
\hskip1em||\endcolouring{}||\endlines
This time syntax-coloured as Python:

\beginlines
\hskip1em||\identifiersyntaxcolouring{}|factorial|\plainsyntaxcolouring{}| = |\constantsyntaxcolouring{}|1|\plainsyntaxcolouring{}||
\hskip1em||\reservedsyntaxcolouring{}|for|\plainsyntaxcolouring{}| |\identifiersyntaxcolouring{}|i|\plainsyntaxcolouring{}| |\reservedsyntaxcolouring{}|in|\plainsyntaxcolouring{}| |\identifiersyntaxcolouring{}|range|\plainsyntaxcolouring{}|(|\constantsyntaxcolouring{}|2|\plainsyntaxcolouring{}|, |\identifiersyntaxcolouring{}|n|\plainsyntaxcolouring{}| + |\constantsyntaxcolouring{}|1|\plainsyntaxcolouring{}|):|
\hskip1em|    |\identifiersyntaxcolouring{}|factorial|\plainsyntaxcolouring{}| *= |\identifiersyntaxcolouring{}|i|\plainsyntaxcolouring{}||
\hskip1em||
\hskip1em||\identifiersyntaxcolouring{}|print|\plainsyntaxcolouring{}|(|\identifiersyntaxcolouring{}|factorial|\plainsyntaxcolouring{}|)|
\hskip1em||\endcolouring{}||\endlines

\weavesections{10}{Images and gadgets}{Chapter 2: Commentary - Markdown Demonstration\quad$\S$10}{\S}{2/md}%
\inwebanchor{para2_md10}\inwebimagewidth{inweb/Tests/omnium/Figures/stlucia.jpg}{15}















\weavesections{11}{Carousels}{Chapter 2: Commentary - Markdown Demonstration\quad$\S$11}{\S}{2/md}%
\inwebanchor{para2_md11}So this is a carousel by Markdown:


\medskip\hrule\smallskip
\beginlines
\hskip1em||\plainsyntaxcolouring{}|    |\functionsyntaxcolouring{}|ROOT|\plainsyntaxcolouring{}| ---> |\functionsyntaxcolouring{}|DOCUMENT|\plainsyntaxcolouring{}||
\hskip1em||\endcolouring{}||\endlines
\centerline{\bf 1/3. Stage 1 - Raw tree}


\smallskip\hrule\medskip

\medskip\hrule\smallskip
\beginlines
\hskip1em||\plainsyntaxcolouring{}|    |\functionsyntaxcolouring{}|ROOT|\plainsyntaxcolouring{}| ---> |\functionsyntaxcolouring{}|DOCUMENT|\plainsyntaxcolouring{}||
\hskip1em|                |{\tt\|}||
\hskip1em|              |\functionsyntaxcolouring{}|NODE|\plainsyntaxcolouring{}| |\constantsyntaxcolouring{}|1|\plainsyntaxcolouring{}|  ---  |\functionsyntaxcolouring{}|NODE|\plainsyntaxcolouring{}| |\constantsyntaxcolouring{}|2|\plainsyntaxcolouring{}|  ---  |\functionsyntaxcolouring{}|NODE|\plainsyntaxcolouring{}| |\constantsyntaxcolouring{}|3|\plainsyntaxcolouring{}|  --- ...|
\hskip1em||\endcolouring{}||\endlines
\centerline{\bf 2/3. Stage 2 - Developed tree}


\smallskip\hrule\medskip

\medskip\hrule\smallskip
\beginlines
\hskip1em||\plainsyntaxcolouring{}|    |\functionsyntaxcolouring{}|ROOT|\plainsyntaxcolouring{}| ---> |\functionsyntaxcolouring{}|DOCUMENT|\plainsyntaxcolouring{}||
\hskip1em|                |{\tt\|}||
\hskip1em|              |\functionsyntaxcolouring{}|NODE|\plainsyntaxcolouring{}| |\constantsyntaxcolouring{}|1|\plainsyntaxcolouring{}|  ---  |\functionsyntaxcolouring{}|NODE|\plainsyntaxcolouring{}| |\constantsyntaxcolouring{}|2|\plainsyntaxcolouring{}|  ---  |\functionsyntaxcolouring{}|NODE|\plainsyntaxcolouring{}| |\constantsyntaxcolouring{}|3|\plainsyntaxcolouring{}|  --- ...|
\hskip1em|                |{\tt\|}|            |{\tt\|}|            |{\tt\|}||
\hskip1em|              |\elementsyntaxcolouring{}|text|\plainsyntaxcolouring{}| |\constantsyntaxcolouring{}|1|\plainsyntaxcolouring{}|       |\elementsyntaxcolouring{}|text|\plainsyntaxcolouring{}| |\constantsyntaxcolouring{}|2|\plainsyntaxcolouring{}|       |\elementsyntaxcolouring{}|text|\plainsyntaxcolouring{}| |\constantsyntaxcolouring{}|3|\plainsyntaxcolouring{}|  ...|
\hskip1em||\endcolouring{}||\endlines
\centerline{\bf 3/3. Stage 3 - Completed tree}


\smallskip\hrule\medskip

\weavesections{12}{A gratuitously long table}{Chapter 2: Commentary - Markdown Demonstration\quad$\S$12}{\S}{2/md}%
\inwebanchor{para2_md12}
\medskip\centerline{\vbox{\offinterlineskip
\hrule
\halign{&\vrule#&\strut\quad\hfil#\hfil\quad&\vrule#&\strut\quad\hfil#\hfil\quad&\vrule#&\strut\quad\hfil#\hfil\quad\cr
height2pt&\omit&&\omit&&\omit&\cr
&Atomic number&&Symbol&&Name&\cr
height2pt&\omit&&\omit&&\omit&\cr
\noalign{\hrule}
height2pt&\omit&&\omit&&\omit&\cr
&1&&H&&Hydrogen&\cr
&2&&He&&Helium&\cr
&3&&Li&&Lithium&\cr
&4&&Be&&Beryllium&\cr
&5&&B&&Boron&\cr
&6&&C&&Carbon&\cr
&7&&N&&Nitrogen&\cr
&8&&O&&Oxygen&\cr
&9&&F&&Fluorine&\cr
&10&&Ne&&Neon&\cr
&11&&Na&&Sodium&\cr
&12&&Mg&&Magnesium&\cr
&13&&Al&&Aluminium&\cr
&14&&Si&&Silicon&\cr
&15&&P&&Phosphorus&\cr
&16&&S&&Sulfur&\cr
&17&&Cl&&Chlorine&\cr
&18&&Ar&&Argon&\cr
&19&&K&&Potassium&\cr
&20&&Ca&&Calcium&\cr
&21&&Sc&&Scandium&\cr
&22&&Ti&&Titanium&\cr
&23&&V&&Vanadium&\cr
&24&&Cr&&Chromium&\cr
&25&&Mn&&Manganese&\cr
&26&&Fe&&Iron&\cr
&27&&Co&&Cobalt&\cr
&28&&Ni&&Nickel&\cr
&29&&Cu&&Copper&\cr
&30&&Zn&&Zinc&\cr
&31&&Ga&&Gallium&\cr
&32&&Ge&&Germanium&\cr
&33&&As&&Arsenic&\cr
&34&&Se&&Selenium&\cr
&35&&Br&&Bromine&\cr
&36&&Kr&&Krypton&\cr
&37&&Rb&&Rubidium&\cr
&38&&Sr&&Strontium&\cr
&39&&Y&&Yttrium&\cr
&40&&Zr&&Zirconium&\cr
&41&&Nb&&Niobium&\cr
&42&&Mo&&Molybdenum&\cr
&43&&Tc&&Technetium&\cr
&44&&Ru&&Ruthenium&\cr
&45&&Rh&&Rhodium&\cr
&46&&Pd&&Palladium&\cr
&47&&Ag&&Silver&\cr
&48&&Cd&&Cadmium&\cr
&49&&In&&Indium&\cr
&50&&Sn&&Tin&\cr
&51&&Sb&&Antimony&\cr
&52&&Te&&Tellurium&\cr
&53&&I&&Iodine&\cr
&54&&Xe&&Xenon&\cr
&55&&Cs&&Caesium&\cr
&56&&Ba&&Barium&\cr
&57&&La&&Lanthanum&\cr
&58&&Ce&&Cerium&\cr
&59&&Pr&&Praseodymium&\cr
&60&&Nd&&Neodymium&\cr
&61&&Pm&&Promethium&\cr
&62&&Sm&&Samarium&\cr
&63&&Eu&&Europium&\cr
&64&&Gd&&Gadolinium&\cr
&65&&Tb&&Terbium&\cr
&66&&Dy&&Dysprosium&\cr
&67&&Ho&&Holmium&\cr
&68&&Er&&Erbium&\cr
&69&&Tm&&Thulium&\cr
&70&&Yb&&Ytterbium&\cr
&71&&Lu&&Lutetium&\cr
&72&&Hf&&Hafnium&\cr
&73&&Ta&&Tantalum&\cr
&74&&W&&Tungsten&\cr
&75&&Re&&Rhenium&\cr
&76&&Os&&Osmium&\cr
&77&&Ir&&Iridium&\cr
&78&&Pt&&Platinum&\cr
&79&&Au&&Gold&\cr
&80&&Hg&&Mercury&\cr
&81&&Tl&&Thallium&\cr
&82&&Pb&&Lead&\cr
&83&&Bi&&Bismuth&\cr
&84&&Po&&Polonium&\cr
&85&&At&&Astatine&\cr
&86&&Rn&&Radon&\cr
&87&&Fr&&Francium&\cr
&88&&Ra&&Radium&\cr
&89&&Ac&&Actinium&\cr
&90&&Th&&Thorium&\cr
&91&&Pa&&Protactinium&\cr
&92&&U&&Uranium&\cr
&93&&Np&&Neptunium&\cr
&94&&Pu&&Plutonium&\cr
&95&&Am&&Americium&\cr
&96&&Cm&&Curium&\cr
&97&&Bk&&Berkelium&\cr
&98&&Cf&&Californium&\cr
&99&&Es&&Einsteinium&\cr
&100&&Fm&&Fermium&\cr
&101&&Md&&Mendelevium&\cr
&102&&No&&Nobelium&\cr
&103&&Lr&&Lawrencium&\cr
&104&&Rf&&Rutherfordium&\cr
&105&&Db&&Dubnium&\cr
&106&&Sg&&Seaborgium&\cr
&107&&Bh&&Bohrium&\cr
&108&&Hs&&Hassium&\cr
&109&&Mt&&Meitnerium&\cr
&110&&Ds&&Darmstadtium&\cr
&111&&Rg&&Roentgenium&\cr
&112&&Cn&&Copernicium&\cr
&113&&Nh&&Nihonium&\cr
&114&&Fl&&Flerovium&\cr
&115&&Mc&&Moscovium&\cr
&116&&Lv&&Livermorium&\cr
&117&&Ts&&Tennessine&\cr
&118&&Og&&Oganesson&\cr
height2pt&\omit&&\omit&&\omit&\cr
}\hrule}}\medskip


\weavesections{13}{Index}{Chapter 2: Commentary - Markdown Demonstration\quad$\S$13}{\S}{2/md}%
\inwebanchor{para2_md13}This sentence should be followed by a mechanically-constructed index.



\beginlines
functions
\quad |collatz|\hskip1em {\bf {\inweblinktoanchor{para1_ld4}{1/ld:$\S$4}}}
\quad |lothar|\hskip1em {\bf {\inweblinktoanchor{para1_ld6}{1/ld:$\S$6}}}
index\hskip1em {\inweblinktoanchor{para2_md1}{$\S$1}}, {\bf {\inweblinktoanchor{para2_md5}{$\S$5}}}
/non-standard/\hskip1em {\inweblinktoanchor{para2_md1}{$\S$1}}, {\bf {\inweblinktoanchor{para2_md5}{$\S$5}}}
|Z_X_1|\hskip1em {\inweblinktoanchor{para2_md1}{$\S$1}}, {\bf {\inweblinktoanchor{para2_md5}{$\S$5}}}
\endlines

% End of weave
\end

